%% ma: L10-B13-0047
%% lop: 10
%% chuong: 5
%% bai: 13
%% dang: 10.13.8
%% loai: DS
%% muc: [NB, NB, TH, TH]
%% dap_an: "SĐSS"
%% nguon: "chuyen-de-cac-so-dac-trung-cua-mau-so-lieu-khong-ghep-nhom-toan-10-chuong-trinh-moi.pdf (GV. Phan Nhật Linh) (P02), Bài 02 Dạng 1 Phần II Câu 2"
%% kien_thuc: "Lớp 10 Bài 13 (số trung bình, trung vị)"
%% trang_thai: chinh-thuc
%% ghi_chu: "Tự duyệt theo yêu cầu giáo viên 22/10/2026."
\begin{ex}
Số liệu sau cho biết lãi (đã quy tròn) hàng tháng của một cửa hàng trong năm 2022 (đơn vị: triệu đồng):
\begin{center}
\begin{tabular}{|c|c|c|c|c|c|c|c|c|c|c|c|c|}\hline
Tháng & $1$ & $2$ & $3$ & $4$ & $5$ & $6$ & $7$ & $8$ & $9$ & $10$ & $11$ & $12$\\\hline
Lãi & $12$ & $15$ & $18$ & $13$ & $13$ & $16$ & $18$ & $14$ & $15$ & $17$ & $20$ & $17$\\\hline
\end{tabular}
\end{center}
Xét tính đúng sai của các khẳng định sau:
\choiceTF{Lãi thấp nhất của cửa hàng là $13$ triệu đồng}{\True Sắp xếp các số trong mẫu theo thứ tự không giảm được: $12$; $13$; $13$; $14$; $15$; $15$; $16$; $17$; $17$; $18$; $18$; $20$}{Số trung bình của mẫu là $\overline x\approx13{,}67$ (triệu đồng)}{Trung vị của mẫu là $16$}
\loigiai{a) Sai: lãi thấp nhất là $12$ triệu đồng.

b) Đúng.

c) Sai: $\overline x=\dfrac{188}{12}\approx15{,}67$.

d) Sai: $M_e=\dfrac{15+16}2=15{,}5$.}
\end{ex}
